% !PRINTED 566
\weberpage{566}

\begin{center}
§ 151. \textbf{The Complex Multiplication of the Function} $\wp(u)$.
\end{center}

We now assume that the periods $\omega_1$, $\omega_2$ of the function $\wp(u)$ satisfy a quadratic equation
\weq{1}{A\omega_2^2 + B\omega_2\omega_1 + C\omega_1^2 = 0,}
in which $A$, $B$, $C$ are integral rational numbers without a common divisor. Let $A$ be positive, and let the discriminant of the quadratic form $(A, B, C)$
\weq{2}{B^2 - 4AC = \Delta}
be a \gesperrt{negative fundamental discriminant}, i.\,e. let $\Delta$ contain no square factor after whose separation a discriminant remains (§ 84).

If we set
\weq{3}{\omega = \dfrac{\omega_2}{\omega_1},}
then from (1) there follows:
\weq{4}{2A\omega = -B + \sqrt{\Delta},}
and if we assume $\sqrt{\Delta}$ to be positive imaginary, then $\omega$ receives a positive imaginary part and can serve as the modulus of a $\vartheta$-function. $\Delta$ is the fundamental number of an imaginary quadratic field, which we denote by $\Omega$\textsuperscript{1)}\begingroup\renewcommand{\thefootnote}{1)}\footnotetext[1]{The assumption that $\Delta$ is a fundamental discriminant is made here for the sake of simplification. In my memoir ``Über Zahlgruppen in algebraischen Körpern'' (Mathematische Annalen, Bd. L) this assumption is not made. One then obtains more general fields, which belong to the orders, as in § 124.}\endgroup.

The class field $\mathfrak{R}(\Delta)$ is determined by the singular invariant $j(\omega)$, and its relative degree with respect to $\Omega$ is equal to the class number of the discriminant $\Delta$. The integers of the quadratic field $\Omega = \mathfrak{R}(\sqrt{\Delta})$ are of the form:
\weq{5}{\mu = \dfrac{x + y\sqrt{\Delta}}{2},}
in which $x$ and $y$ are integral rational numbers that satisfy the condition
\[ x \equiv By \pmod{2}. \]
For the number $\mu$ defined by (5) we determine four integral rational numbers $a$, $b$, $c$, $\partial$:
\weq{6}{\begin{aligned} a &= \dfrac{x+By}{2}, & b &= Ay, \\ c &= -Cy, & \partial &= \dfrac{x-By}{2}, \end{aligned}}
%CONT>>
